\chapter{Discussion}
\label{Discussion}

This chapter interprets the evaluation results (Section~\ref{sec:disc-interp}), analyses the role of union in multimodal design (Section~\ref{sec:disc-union}), acknowledges threats to validity (Section~\ref{sec:disc-threats}), and summarises the contributions (Section~\ref{sec:disc-contrib}).

\section{Interpretation of Results}
\label{sec:disc-interp}

\textbf{RQ1: Object re-identification.} The result is unequivocal and the effect size is maximal. The re-ID mechanism transforms visual detection from zero to an F1 of up to 0.839. All four VLM models, spanning two different architectures (Mistral-based Pixtral/Ministral and Google-based Gemini), fail identically without re-ID. This generality is significant: it demonstrates that the zero-event bug is a fundamental limitation of VLM-based interval detection, not an artefact of a particular model or provider. The re-ID mechanism is simple ,  two-phase prompt engineering with block-overlap matching ,  yet it is sufficient to enable meaningful event detection. No separate tracking model or training data is required, which is important for reproducibility and for application to future VLM architectures.

\textbf{RQ2: Audio model ablation.} Qwen2-Audio achieves perfect precision (1.000) for all window configurations of 2.5 s or larger. Every detection is correct. This is a critical property for the multimodal union: false positives from audio would propagate to condition C and degrade its precision. The LALM's precision is not matched by PANNs CNN14, whose highest F1 (0.429) is less than 60\% of the LALM's best micro F1 (0.788). The window/hop ablation reveals that 2.5 s windows with 1.25 s hop are optimal for both providers: this balances acoustic context (the model sees enough signal to identify the sound class) with temporal resolution (events are localized precisely enough for interval queries).

The per-event analysis reveals distinct model strengths. Qwen2-Audio is strongest on vehicle escape (perfect recall, 5/5) and gunshot (5/5), but fails on fight detection entirely (0/5). The fight failure is instructive: the shout-to-impact temporal pattern may be too subtle or too tightly spaced for the 2.5 s window to capture both sounds as distinct events, or the fight sounds in the synthetic dataset may lack sufficient acoustic distinctiveness. PANNs achieves its best performance on gunshot (5/5) but with lower precision due to false positives, and fails entirely on fight and vehicle escape. This asymmetry suggests that the two models could be complementary in a hierarchical detection design.

\textbf{RQ3: Multimodal vs.\ unimodal.} The union-based condition C achieves 29/30 TP (F1 = 0.892, 6 FP, 1 FN) with Gemini 3.6 Flash + Qwen2-Audio, outperforming both visual-only (best F1 = 0.839) and audio-only (best micro F1 = 0.788). The thesis claim $C \ge \max(A, B)$ holds on every scene. The 1 false negative is a scene where neither modality succeeds, highlighting the remaining gap that stronger perception models could address.

\section{The Role of Union in Multimodal Design}
\label{sec:disc-union}

The design choice of union over temporal cross-modal join merits discussion. The conventional intuition is that multimodal fusion should require coincidence: an audio event is confirmed by a temporally proximate visual event, and vice versa. However, this approach has a fundamental asymmetry in forensic applications. The cost of a false negative (a missed event) far exceeds the cost of a false positive (an unnecessary operator review), because a missed event could mean a missed crime. Temporal join reduces false positives at the expense of increased false negatives ,  the opposite of what a forensic system should optimise for.

Union makes the opposite trade-off: it preserves every detection from every modality, at the cost of potentially more false positives. But the union approach is auditable: every detection carries a modality tag, so an operator reviewing the output can see whether a detection came from audio, visual, or both. The experiment demonstrates that union is the correct operation for forensic surveillance: $C$ consistently dominates $\max(A, B)$ with no exceptions.

\section{Threats to Validity}
\label{sec:disc-threats}

\textbf{Internal validity.} The evaluation uses synthetic videos generated by Dreamina Seedance. Synthetic data allows precise control over event timing and modality signal, but may not capture the variability of real CCTV recordings. The audio tracks are also generated and may have different acoustic properties from real surveillance environments.

\textbf{External validity.} The dataset of 30 scenes is substantially larger than the 8 scenes used in preliminary evaluations of prior work, but it remains limited. The fixed camera angle, single-scene design, and synthetic nature of the footage mean that the absolute F1 values should not be interpreted as deployment-ready accuracy.

\textbf{Construct validity.} The evaluation uses F1 score based on binary verdicts. This captures whether an event was detected at all, but does not measure temporal precision: a detection five seconds offset from the ground truth receives the same credit as a perfectly aligned one. Future work should incorporate temporal IoU (intersection over union) to provide a more granular view of temporal alignment quality.

\textbf{Conclusion validity.} The effect sizes are large: re-ID transforms detection from zero to an F1 of 0.839; Qwen2-Audio outperforms PANNs by 0.359 micro F1 points; union adds 0.053 F1 points over the best unimodal condition (visual F1 = 0.839 to multimodal F1 = 0.892). These magnitudes are large enough that statistical significance is not in question, even with 30 scenes.

\section{Key Contributions}
\label{sec:disc-contrib}

This thesis makes seven contributions:

\begin{enumerate}
  \item \textbf{First LALM for surveillance forensics.} Qwen2-Audio-7B-Instruct achieves perfect precision (1.000) and micro F1 = 0.788, substantially outperforming PANNs CNN14 (F1 = 0.429).
  \item \textbf{Object re-identification for VLM-based surveillance.} A two-phase tracking prompt with block-overlap matching fixes the zero-event bug: all four tested VLM models achieve 0 TP without re-ID.
  \item \textbf{First labeled multimodal forensic dataset.} 30 scenes, 6 event types, per-frame visual + audio + verdict ground truth.
  \item \textbf{Union-based multimodal design.} $C = A \cup B$ preserves all detections while temporal joins reduce recall.
  \item \textbf{VLM benchmark across four models.} Gemini 3.6 Flash is the best performer (F1 = 0.839; 29/30 TP when combined with Qwen2-Audio).
  \item \textbf{Closed-vocabulary LALM prompt.} Forces the LALM to select from predetermined sound classes, eliminating post-hoc normalisation.
  \item \textbf{Auditable forensic pipeline.} Every detection is the output of a readable, editable SQL query ,  no black-box reasoning.
\end{enumerate}
